\documentclass[10pt,a4paper]{amsart}
\usepackage[margin=19mm]{geometry}
\usepackage{amsmath,amssymb,mathtools}
\usepackage[scr=rsfs,cal=euler]{mathalfa}
\usepackage{xfrac}
\usepackage[unicode,hidelinks,bookmarks=false]{hyperref}
\newcommand{\ie}{i.e.\ }
\newcommand{\cf}{cf.\ }
\newcommand{\resp}{respectively\ }
\newcommand{\Subsection}{\subsection}
\providecommand{\nobreakdash}{\nobreak\hbox{-}}
\setlength{\emergencystretch}{2em}
\multlinegap0pt
\usepackage{amssymb}
\usepackage[shortlabels]{enumitem}
\newtheorem{lemma}{Lemma}
\newtheorem{proposition}{Proposition}
\title{On simple transposed Poisson algebras}
\author{Amir Fern\'andez Ouaridi}
\date{Source: arXiv:2604.26115v1 (CC BY 4.0)}
\begin{document}
\maketitle

\noindent\emph{Source and licence.} On simple transposed Poisson algebras. Version arXiv:2604.26115v1 (CC BY 4.0), distributed by arXiv under the Creative Commons Attribution 4.0 International licence, http://creativecommons.org/licenses/by/4.0/, as recorded in the arXiv licence metadata for this exact version and shown on the abstract page https://arxiv.org/abs/2604.26115v1. Full licence notice: \texttt{licenses/arxiv-2604.26115-CC-BY-4.0.txt}.\par\smallskip

\noindent\textit{Editorial scope.} Counted unit: the article's proposition proposition (source0.tex lines 904--907). The article's complete original proof of it is delivered with it (source0.tex lines 909--973, one \texttt{proof} environment of 65 lines). No mathematics is written, repaired or re-derived: the statement and the proof are the article's own lines, in the article's own notation, and no retained line is substituted or re-flowed.  Retained uncounted as the source-local context the proof needs: lines 809--814; lines 867--873.  Nothing was omitted from any retained span, so the package carries no omission record.  No retained line contains a citation, so the wrapper carries no bibliography and no bibliography entry.  No retained line contains a figure environment, an image-inclusion command or drawing code, so no image is delivered and the package declares no asset dependency.  Editorial formatting only, in addition: an AMS \texttt{amsart} preamble that restates the article's own declarations and macros used by the retained lines, namely the environments \texttt{lemma}, \texttt{proposition} on separate counters, together with the standard packages those lines need.  The retained environments print in this document as Lemma 1, Proposition 1 in order of appearance, whereas the article prints the same items with the numbering of its own full document.  The complete native source of the article as distributed in the arXiv e-print for this exact version is bundled unchanged as \texttt{sources/199364/source0.tex}.
\begin{lemma}\label{relsbeta1}
    Let $(\alpha,\beta,V)$ be an irreducible representation of
$\mathcal{W}_n(e_{-1})$. Denote $T=\beta(e_{-1})$. Then for every $x\in \mathcal{W}_n(e_{-1})$, we have
$$\beta(x) = \alpha(x)T - \alpha(\delta(x)) \qquad \textrm{ and } \qquad \beta(x) = T \alpha(x) - 2 \alpha(\delta(x)).$$
Moreover, we have  $\alpha(x)T=T\alpha(x)-\alpha(\delta(x))$. 
\end{lemma}

\begin{lemma}\label{actionalpha}
For every $0\leq i\leq N-2$ and every $k\geq 0$, we have
$$
\alpha(f_i)(v_k)=(-1)^{i+1}\binom{k}{i+1}v_{k-i-1},
$$
where we use the convention $v_s=0$ if $s<0$.
\end{lemma}

\begin{proposition}
Let $(\alpha,\beta,V)$ be an irreducible representation of
$\mathcal{W}_n(e_{-1})$. Then $(\alpha,V)$ is isomorphic to the regular representation.
\end{proposition}

\begin{proof}
Denote $m=\textrm{dim}(V)$. Since $V=V_0$, we may choose the minimal $m$ such that
$V$ is spanned by $\left\{v_0,\dots,v_{m-1}\right\}$. Then $\left\{v_0,\dots,v_{m-1}\right\}$ is a basis of $V$.
We first prove that $N$ divides $m$. By Lemma~\ref{actionalpha}, for every $0\leq i\leq N-2$ we have
$$
\alpha(f_i)(v_m)=(-1)^{i+1}\binom{m}{i+1}v_{m-i-1}.
$$
Since $v_m\in \textrm{span}(v_0,\dots,v_{m-1})$, the left hand side belongs to
$\textrm{span}(v_0,\dots,v_{m-i-2})$. Therefore
$$
\binom{m}{i+1}\equiv 0 \pmod p
$$
for every $0\leq i\leq N-2$. 
Write $m = \sum_{k=0}^{r} m_k p^k$ in base $p$. Then for $0\leq j\leq n-1$ it follows
$$
m_j = \binom{m_j}{1} \equiv \binom{m}{p^j}\equiv 0 \pmod p,
$$
by Lucas' theorem.
We deduce $N\mid m$.
Moreover, we claim $U=\textrm{span}(v_{m-N},v_{m-N+1},\dots,v_{m-1})$ is $\alpha$-stable and $\beta$-stable. 
Indeed, for $m-N\leq k\leq m-1$ and $0\leq i\leq N-2$, Lemma~\ref{actionalpha} gives
$$
\alpha(f_i)(v_k)=(-1)^{i+1}\binom{k}{i+1}v_{k-i-1}.
$$
If $k-i-1\ge m-N$, then $\alpha(f_i)(v_k)\in U$. Otherwise, $k-i-1<m-N$, so
$i+1>k-(m-N).$

Since $N\mid m$ and $N\mid (m-N)$, Lucas' theorem gives
$$
\binom{k}{i+1}
=
\binom{m-N+(k-(m-N))}{i+1}
\equiv
\binom{k-(m-N)}{i+1}
\equiv 0 \pmod p.
$$
Hence, we have $\alpha(f_i)(v_k)\in U$. Thus, $U$ is $\alpha$-stable.
Now, Lemma~\ref{relsbeta1}  and $\alpha(f_i)(v_m) = 0$ yields
$$
\beta(f_i)(v_k)=\alpha(f_i)(Tv_k)-\alpha(\delta(f_i))(v_k)
=\alpha(f_i)(v_{k+1})-\alpha(f_{i-1})(v_k)\in U,
$$
so $U$ is also $\beta$-stable. Hence, the irreducibility of $(\alpha, \beta, V)$ implies $U=V$ and $m=N$.
\medskip

Now, consider the linear map $\varphi:\mathcal{W}_n(e_{-1})\to V$ given by
$
\varphi(x)=\alpha(x)(v_{N-1}).
$
Since $(\alpha,V)$ is a representation of the associative algebra $(\mathcal{W}_n(e_{-1}),\bullet)$, $\varphi$ is a homomorphism between the regular representation and $(\alpha,V)$.
For $0\leq k\leq N-1$, we have
$$
\varphi(f_{N-k-2})
=
\alpha(f_{N-k-2})(v_{N-1})
=
(-1)^{N-k-1}\binom{N-1}{N-k-1}v_k.
$$
Then there is $t\in \mathbb{F}$ such that $\varphi(tf_{N-k-2}) = v_k$, because we have
$$
\binom{N-1}{N-k-1} = \binom{p^n-1}{p^n-k-1}  \not\equiv 0 \pmod p.
$$
Hence $\varphi$ is surjective. Since both spaces have dimension $N$,  the map $\varphi$ is an isomorphism.
Therefore, $(\alpha,V)$ is isomorphic to the regular representation of $(\mathcal{W}_n(e_{-1}),\bullet)$.
\end{proof}

\end{document}
